The first row is any permutation of the three colours, giving $3!=6$ choices. The second row must be a derangement of the first; with three symbols there are exactly two such permutations (the two cyclic shifts). The third row is then forced by the missing colour in each column and automatically contains all three colours. Thus $\boxed{12}$.
\begin{asy}
import graph;import geometry;usepackage("amsmath");size(7cm);
for(int i=0;i<=3;++i){draw((i,0)--(i,3));draw((0,i)--(3,i));}string[] a={"G","W","R"};for(int i=0;i<3;++i)for(int j=0;j<3;++j)label(a[(i+j)%3],(j+.5,2.5-i));label("one of 12 arrangements",(1.5,-.5));
\end{asy}

